\documentclass[11pt]{article}
\usepackage{amsmath,amssymb,amsthm}

\newtheorem{theorem}{Theorem}
\newtheorem{lemma}[theorem]{Lemma}

\begin{document}

\section*{Placement of the end-of-proof mark}

\begin{lemma}\label{lem:sum}
For every $n \ge 1$ we have $1 + 2 + \cdots + n = n(n+1)/2$.
\end{lemma}

\begin{proof}[Proof of Lemma~\ref{lem:sum}]
Write the sum twice, once forwards and once backwards, and add the two rows
term by term. Each of the $n$ columns sums to $n+1$, hence
\[
  2\sum_{k=1}^{n} k = n(n+1). \qedhere
\]
\end{proof}

\begin{theorem}\label{thm:sq}
For every $n \ge 1$, $\sum_{k=1}^{n} (2k-1) = n^2$.
\end{theorem}

\begin{proof}
By Lemma~\ref{lem:sum},
\begin{align*}
  \sum_{k=1}^{n} (2k-1) &= 2\sum_{k=1}^{n} k - n \\
                        &= n(n+1) - n \\
                        &= n^2. \qedhere
\end{align*}
\end{proof}

\begin{theorem}
Let $a, b$ be integers with $a \mid b$ and $b \mid a$. Then $a = \pm b$.
\end{theorem}

\begin{proof}
We treat two cases.
\begin{enumerate}
  \item If $a = 0$ then $b = 0$ as well, since $0$ divides only itself.
  \item If $a \neq 0$, write $b = ka$ and $a = lb$. Then $a = lka$, so
        $lk = 1$ and $k = \pm 1$. \qedhere
\end{enumerate}
\end{proof}

\begin{theorem}
There are infinitely many primes.
\end{theorem}

\begin{proof}
Suppose $p_1, \dots, p_r$ are all the primes and let $N = p_1 \cdots p_r + 1$.
\begin{itemize}
  \item No $p_i$ divides $N$, because it leaves remainder $1$.
  \item But $N > 1$ has a prime factor, which is a contradiction. \qedhere
\end{itemize}
\end{proof}

\begin{proof}[Second proof]
The series $\sum_p 1/p$ diverges, which is impossible for a finite set of
primes. We only sketch this argument here.
\end{proof}

\renewcommand{\qedsymbol}{$\blacksquare$}

\begin{theorem}
The number $\sqrt{2}$ is irrational.
\end{theorem}

\begin{proof}
If $\sqrt{2} = p/q$ in lowest terms, then $p^2 = 2q^2$ is even, so $p$ is
even, so $q$ is even, contradicting lowest terms.
\end{proof}

\begin{proof}[Proof with a displayed equation]
Equivalently, no positive integers $p, q$ satisfy
\begin{equation}
  p^2 = 2q^2 \qedhere
\end{equation}
\end{proof}

\renewcommand{\qedsymbol}{\textsc{qed}}

\begin{proof}[A remark in proof form]
The same argument shows that $\sqrt{m}$ is irrational whenever $m$ is not a
perfect square.
\end{proof}

\end{document}
